% In this paper, we presented \textbf{\appname}, a lightweight localized model editing framework for mitigating package hallucination in LLMs. 
% Instead of treating package hallucination as a standard factual editing problem, \appname formulates it as a package-validity boundary editing problem, where the goal is to improve the model's distinction between valid and hallucinated package names.  
% \appname first identifies modules related to package hallucination through risk-aware localization, and then edits only these localized modules using lightweight LoRA adapters with a boundary-aware objective.
% We evaluated \appname on three representative open-source LLMs: DeepSeekCoder, Qwen3, and Llama-3.1. 
% The results show that \appname effectively reduces package hallucination in package recommendation, lowering Package-HR by 79.9\% on edit prompts and by 65.4\% on unseen prompts. 
% The edited package-validity boundary also transfers to other package-related tasks, reducing Package-HR by 12.8\% in code generation and by 34.0\% in \texttt{pip install} recommendation.
% Compared with existing baselines, \appname better balances hallucination reduction and valid package preservation, while several baselines reduce hallucinations by suppressing package generation.
% Our ablation study further confirms the importance of each component in \appname, and the cost analysis shows that \appname is practical and lightweight.  
% Overall, \appname provides a practical step toward improving the reliability of LLM-generated package recommendations without full model retraining or external retrieval pipelines. 
% Future work will extend \appname to more package ecosystems and broader dependency management scenarios.

In this paper, we presented \textbf{\appname}, a lightweight localized model editing framework for mitigating package hallucination in LLMs. 
\appname formulates package hallucination as a package-validity boundary editing problem, where the goal is to improve the model's distinction between valid and hallucinated package names. 
It first locates modules related to package hallucination through risk-aware localization, and then edits only these modules with lightweight LoRA adapters and a boundary-aware objective.
We evaluated \appname on three representative open-source LLMs: DeepSeekCoder, Qwen3, and Llama-3.1. 
The results show that \appname effectively reduces package hallucination in package recommendation, lowering Package-HR by 79.9\% on edit prompts and by 65.4\% on unseen prompts. 
The edited package-validity boundary also transfers to other package-related tasks, reducing Package-HR by 12.8\% in code generation and by 34.0\% in \texttt{pip install} recommendation. 
Compared with existing baselines, \appname better balances hallucination reduction and valid package preservation, while several baselines reduce hallucinations by suppressing package generation.
Ablation and cost analyses further show that each component contributes to \appname's effectiveness, and the framework is practical and lightweight. 
Overall, \appname provides a practical step toward mitigating package hallucinations in LLMs without full model retraining or external retrieval pipelines. 
Future work will extend \appname to more package ecosystems and broader dependency management scenarios.